\section{Introduction}\label{sec:introduction} % ! SHOULDN'T BE HUGE, 1-2 pages

% * ____ proposed paragraphs ________________
% ! Specify interest to this conference, how my papers are connected to the topics selected for the conference

% ! Identify a gap between existing tools and my solution, WHAT IS THE DIFFERENCE
Nowadays climate change is one of the most popular questions all around the world. Accordingly to the Peoples' 
Climate Vote 2024~\cite{undp2024pcv}, approximately 7 out of 10 people in the world says that climate 
change influences big decisions in their families.
Understanding the tendency where climate is moving, temperature changes, usual precipitation, solar radiation 
in specific region, ability to compare climate between different periods of time or between cities
could be crutial to make conclusions about future, take correct decisions in agriculture, tourism, and other areas.
However, access to climate data is not that easy for everyone. We have a lot of climate data yet the big question here:
"How to use it quickly, easily and without any special knowledge?".
Wordlclim is one of the open datasets, with global coverage which provides monthly climate data in rasters between 1970 and 2024, 
with variables such as temperature, precipitation, solar radiation and wind speed, among others (26 variables in total), 
different resolutions of rasters (from 30 arc-seconds to 10 arc-minutes) and different time periods (historical, current and future).
The problem here is that despite the rich data provided by WorldClim, it is combersome to retrieve and visualize 
needed data, it takes a lot of time and requires special knowladge and special GIS software to work with it. So that, 
big quantity of potential users are not able to use it. 

% ? context: worldclim, climate data, how it is used and who use it
% * climate data is important... there are a lot of dataset and worldclim is..
% * worldclim could be used in a lot of cases but it was hard to access this data
% \subsection{Motivation and Background}\label{sec:motivation}
% \cite{S.2026Design}

% ? problem to solve: existing tools are not user friendly, require GIS knowledge, or are static..
% \subsection{Objectives}\label{sec:objectives}
Meanwhile, enternet is full of different web applications which provide climate data visualisation, but they are 
either static and not interactive or require special GIS knowledge to use them. they cover only areas researched
from metheorological stations, which means that for a lot of areas its just impossible to check climate data and 
visualize it. What is more, some of them are not free and require registration. 
At the same time, worldclim doesnt have such limitations as narrow regional coverage,
coarse or point-based resolution, restricted time periods. Indeed, it is a gem, which needs to be polished and 
made accessible for everyone

% ? proposed solution
And to solve this gap between rich data and required knowledge to use it I propose a web application which provides interactive access 
to WorldClim data, with ability to easily visualize neaded period of time, compare it and export results for neaded researches.

% ? structure of the paper